Las soluciones basadas en APIs comerciales como las de OpenAI o Anthropic ofrecen una gran potencia con una inversión inicial casi nula, pero introducen riesgos importantes en lo que respecta a la privacidad y los costes a largo plazo. Uno de los principales inconvenientes es que el envío de datos sensibles de los alumnos a servidores externos puede comprometer la seguridad de la información y dificultar el cumplimiento de normativas estrictas como el GDPR.

Las APIs funcionan mediante el pago por tokens, y esto suele disparar los costes operativos de forma impredecible cuando hay muchos usuarios conectados simultáneamente, lo cual resulta inasumible en entornos educativos con presupuestos limitados. Aunque montar un servidor propio requiere comprar \textit{hardware} de antemano, esta opción ofrece un modelo de costes predecibles, la inversión inicial es significativamente más alta, pero controlable. Investigaciones recientes sitúan el umbral de rentabilidad o \textit{break-even point} en un volumen aproximado de entre 1,5 y 2 millones de tokens diarios, ya que a partir de ese nivel de uso la infraestructura propia compensa el gasto acumulado en cuotas de suscripción externa \cite{pan2025cost}. Alcanzar estas cifras no es raro en el ámbito académico, incluso sin un uso intensivo. Esta estrategia permite rentabilizar la inversión inicial y al mismo tiempo asegura la soberanía tecnológica de la institución al no depender de proveedores externos que puedan modificar de manera unilateral sus términos de uso o condiciones de servicio.

Si bien es cierto que igualar la capacidad de razonamiento de un modelo comercial masivo mediante una infraestructura local es una tarea compleja, un tutor socrático orientado a estudiantes principiantes de ingeniería informática no requiere tal nivel de sofisticación. En este escenario, el diseño de la plataforma debe priorizar la privacidad y el control sobre los datos del alumnado por encima de una potencia de procesamiento que no llegaría a ser aprovechada en su totalidad.